\section{Deformations at cone points with zero projection}
\label{sec:rigidity}

Fix a component $S_p$ of the local deformation base at each cone point
$p$, and let $B=\Def^S(X)$ be the inverse image of their product under
the local classifying maps. We use the scheme structure on this inverse
image. Write $F$ for the set of cone points at which the projection of
$K_S$ to the selected local summand is zero. We first study deformations
of the vertex blow-up at a point of $F$.

\begin{lemma}\label{lem:blowup-rigidity}
Let $U$ be a sufficiently small Stein representative of the
anticanonical cone over a smooth del Pezzo surface $E$ of degree six.
Let $f:V\to U$ be the blow-up of its vertex. Then
\[
 f_*\Theta_V=\Theta_U,\qquad R^1 f_*\Theta_V=0.
\]
Every marked deformation of $V$ over a local Artin $\CC$-algebra is
trivial. Across a small extension the trivialization can be chosen to
extend any specified trivialization of the preceding order.
\end{lemma}

\begin{proof}
The singular locus of $U$ consists of the vertex. The local flow of a
vector field on $U$ therefore preserves the vertex ideal and lifts to
its blow-up. In the opposite direction a vector field descends on the
punctured germ and extends across the vertex by reflexivity of
$\Theta_U$. This proves the direct-image assertion.

Near the exceptional divisor $V$ is a neighborhood of the zero section in
the total space of $K_E$. If $I$ is its ideal, the graded quotients of
$\Theta_V$ for the $I$-adic filtration are
\[
 \bigl(\Theta_E\oplus K_E\bigr)\otimes\mathcal O_E(-kK_E),
 \qquad k\geq 0.
\]
Their first cohomology vanishes, as in the local calculation
in~\cite[proof of Lemma~5.2, (47)--(48)]{W5}. Indeed,
$H^1(E,\mathcal O_E((1-k)K_E))=0$ by rationality, Serre duality and
Kodaira vanishing. For the tangent summand, the toric Euler sequence
has middle term $\bigoplus_{j=1}^6\mathcal O_E(C_j)$, where the $C_j$
are the six boundary curves, and left term $\mathcal O_E^{\oplus4}$.
When $k=0$, the restriction $\mathcal O_{C_j}(C_j)$ has degree $-1$,
giving $H^1(E,\mathcal O_E(C_j))=0$. When $k\geq1$, the divisor
$C_j-kK_E$ has nonnegative degree on each of the six generators of the
cone of curves, so toric vanishing applies. The last term needed in the
cohomology sequence is $H^2(E,\mathcal O_E(-kK_E))=0$.
Induction on the thickenings of $E$ and the theorem on formal functions
give $R^1f_*\Theta_V=0$. Here the powers of $I$ are cofinal with the
pullbacks of the powers of the vertex ideal.

Cartan's theorem~B and Leray now give $H^1(V,\Theta_V)=0$. Consider a
small extension of local Artin algebras with kernel $J$. Once a
trivialization over its quotient is fixed, the differences of local
trivializations lifting it define a class in
$H^1(V,\Theta_V)\otimes J$. The local lifts exist since $V$ is smooth.
The class vanishes, so the lifts can be made to agree. Induction from
the given central marking proves the assertion.
\end{proof}

We shall use the following consequence of semiuniversality. Let $(G,0)$
be a semiuniversal base of an isolated analytic singularity $(U,p)$,
and let $D$ be a formal power-series ring over $\CC$. If a map
$\operatorname{Spf}D\to\widehat G$ induces a trivial marked local
deformation at every Artin quotient, that map is constant. To see this,
choose a generic formal line through the origin on which a first nonzero
term survives, and call its order $m$. The map modulo $t^{m+1}$ factors
through the first infinitesimal neighborhood of $0\in G$, with tangent
coefficient in $(t^m)/(t^{m+1})$. The Kodaira--Spencer isomorphism
identifies this coefficient with a nonzero element of
$T^1_{U,p}\otimes (t^m)/(t^{m+1})$, using the product trivialization
modulo $t^m$. A change of lower-order trivialization cannot erase this
class. In fact every marked automorphism of the product modulo $t^m$
lifts to the product modulo $t^{m+1}$. Its logarithm is a derivation
with coefficients in $(t)$: the finite binomial formula for powers of a
unipotent automorphism is multiplicative at every integer exponent,
hence at every complex exponent, and differentiation at zero gives
the Leibniz rule. Lift the coefficients of this derivation and
exponentiate. Any automorphism of the central germ lifts constantly.
Thus a trivial deformation modulo $t^{m+1}$ would have zero class for
the specified product identification, a contradiction.

\begin{theorem}\label{thm:regular}
For every choice of the local components $S_p$, the analytic germ
$\Def^S(X)$ is smooth of dimension
\[
 h^1(X,\Theta_X)+\dim_{\QQ} K_S.
\]
At every $p\in F$ the local classifying map is constant on this entire
germ. Let $X_F\to X$ be the blow-up of the vertices in $F$, and let
$S'$ denote the induced choices at the remaining cone points.
Blow-down induces an isomorphism of the completions of
$\Def^{S'}(X_F)$ and $\Def^S(X)$.
\end{theorem}

\begin{proof}
The blow-up $\pi:X_F\to X$ is projective and crepant. It replaces each
vertex in $F$ by the smooth zero section of its canonical bundle and
is an isomorphism near the remaining singularities. Rationality of
these cone singularities gives $R\pi_*\mathcal O_{X_F}=\mathcal O_X$.
In particular $X_F$ is projective with trivial dualizing sheaf and
$H^1(X_F,\mathcal O_{X_F})=0$. The same smooth threefold $Y$ resolves
$X$ and $X_F$. Its marked curves at the remaining singularities are
unchanged. Extension of a tuple by zero at $F$ identifies
$K_{S'}(X_F)$ with $K_S$, since every relation in $K_S$ already
vanishes there. In particular every remaining selected cone summand
has nonzero projection.

Write $B'=\Def^{S'}(X_F)$. The selected-base theorem
\cite[Theorem~4.5]{W5} makes $B'$ smooth. Let
$D=\widehat{\mathcal O}_{B',0}$. Analytic blow-down over Artin
algebras gives a formal morphism from $\operatorname{Spf}D$ to
$\widehat{\Def}(X)$ compatible with restriction to open sets and with
base change \cite[Lemmas~4.2--4.3]{W5}. At each vertex in $F$,
Lemma~\ref{lem:blowup-rigidity} trivializes its source over the
inverse image of a Stein neighborhood, to every order. Pushing forward
the product structure sheaf gives the trivial deformation of that
neighborhood. The preceding consequence of semiuniversality shows
that all coordinates of its fixed local classifying map therefore
vanish in $D$.

At every other cone the morphism $\pi$ is a local isomorphism. The
differential of blow-down on its local components agrees with the
original restriction, using the pullback of one global volume form
for the homological identifications. The selected tangent sequence
\cite[Lemma~4.4]{W5} for $X_F$ and the nonzero projection at this cone
show that at least one coordinate of the chosen component has nonzero
linear term in $D$. The completed local deformation ring of the cone is
$\CC[[s,u,v]]/(su,sv)$, with line $(u,v)=0$ and plane $s=0$.
Since $D$ is a domain, the complementary coordinates vanish. Thus the
formal morphism factors through $\widehat B$ at every cone. Nodes
impose no further choice. Denote the factored morphism by $\beta$.

The map $d\beta$ identifies the two selected tangent spaces. Indeed,
Lemma~\ref{lem:blowup-rigidity} and the local isomorphism away from $F$
give $\pi_*\Theta_{X_F}=\Theta_X$ and $R^1\pi_*\Theta_{X_F}=0$.
Leray identifies $H^1(X_F,\Theta_{X_F})$ with $H^1(X,\Theta_X)$.
On locally trivial deformations this is the differential of blow-down:
a vector-field cocycle on overlaps descends by the direct image of
its vector fields. On the local quotient the differential extends the
tuple by zero at $F$, hence gives the identification of relation
spaces above. For each of the pairs $(T,S_T)=(X_F,S')$ and $(X,S)$
we have the exact sequence
\[
 0\longrightarrow H^1(T,\Theta_T)
 \longrightarrow T_0\Def^{S_T}(T)
 \longrightarrow K_{S_T}(T)\otimes_{\QQ}\CC\longrightarrow0.
\]
The isomorphisms on its two ends give the assertion about $d\beta$.

Write $\widehat{\mathcal O}_{B,0}=\CC[[x_1,\ldots,x_r]]/I$
minimally. The images of the $x_i$ in the regular complete local ring
$D$ form a basis of its cotangent space. They form a system of formal
coordinates, so the composite $\CC[[x_1,\ldots,x_r]]\to D$ is an
isomorphism. It follows that $I=0$ and that $\beta$ is an isomorphism.
Regularity of the completion proves smoothness of $B$, and the
tangent sequence gives its dimension. At any $p\in F$ the local
classifying coordinates vanish after pullback to $D$. They therefore
vanish in $\widehat{\mathcal O}_{B,0}$ and in its analytic local
ring. This proves constancy on $B$.
\end{proof}
